\section{보안 관점에서 정확히 무엇을 얻었는가}
\label{sec:security-boundary}

\subsection{국소 무손실(zero loss)의 의미}

\cref{thm:challenge-suffix}는 확률분포 사이의 근사나 random-oracle 이상화를
사용하지 않는다. 같은 전제 아래 두 결정론적 생성 절차의 반환 상태가 같다는
관계적 동치다. 따라서 이 특정 seam에 대해서만
\[
  \delta_{\mu,\mathrm{raw\_top}}=0
\]
이라고 기록할 수 있다.

\begin{boundary}[승격할 수 없는 등식]
\label{bdry:no-promotion}
국소 등식은
\[
  \delta_{\mathrm{Sign}}=0,
  \qquad
  \delta_{\mathrm{Verify}}=0,
  \qquad
  \delta_{\mathrm{Encoding}}=0
\]
중 어느 것도 함의하지 않는다. 검증 스크립트는 이러한 부당한 승격 문자열이
작성 트리에 나타나면 실패하도록 되어 있다.
\end{boundary}

또한 \cref{thm:accepted-api-adapter}의 이름에 \identifier{zero_loss}가 들어가도
현재 허용되는 API 등식은 아니다. 교차 호출 프레임(cross-call frame), 구간
국소 해시 관계와 정확 순차 추적은 이제 컴파일되지만, 생성 해시의 반환값을
완료된 추적의 저장된 관찰값으로 운반하는 생성/재전달 결론이 없으므로
\[
  \delta_{\mu,\mathrm{raw\_api\_accept}}=0
\]
은 아직 기록하지 않는다.

이 구분은 \cref{eq:implementation-bound}의 손실 항을 과소평가하지 않기 위해
필수적이다. 현재 정리는 \(\delta_{\mathrm{Sign}}\)와
\(\delta_{\mathrm{Verify}}\) 내부의 한 결정론적 부분을 0으로 만들었을 뿐,
sampler, rejection, packing, parsing, public context, API memory를 포함한 전체
하이브리드를 닫지 않았다. 다만 API memory의 많은 하위 경계---주소 왕복,
KeyGen 순차 export, Sign trace와 output frame, accepted Verify tail과 입력
바인딩, region-local hash read relation---는 이제 실제 caller 정리로 닫혔다.

Week~8--10의 HBZ 잎/표/단일 단계, 순수 바이트 스택, 실제 접미부 복사와 결합
제어, 그리고 실제 인코더의 내부 바이트/프레임 단계와 성공 시 크기 범위도
결정론적 기능정확성 증거다. Week~11의 전체 실제 인코더 성공 접미부 등식까지
같은 종류의 증거로 닫혔다. Week~12의 실제 디코더 정확 추적 복원도
정확한 입력 관계 아래의 결정론적 부분정확성으로 닫혔다. Week~13은 세 실제
절차의 순차 결합까지 성공 조건부 부분정확성으로 닫았다. Week~14는 실제
prepare/apply를 포함한 full-HBZ 직접 결합과 production SignaturePack/Unpack
경계의 정확 동치까지 닫았다. Week~15는 고정 all-zero 입력에서 종료한
production 실행의 zero-size 실패를 배제했다. 그러나 이 고정 입력의 종료와
losslessness, 일반 정준 입력 성공, (h) 코덱, metadata/padding 및 전체 parse가
열려 있으므로 이를 \(\delta_{\mathrm{Encoding}}=0\)으로 승격할 수 없다.

Week~16의 KeyGen 스냅샷 정리도 결정론적 기능정확성 증거지만
\(\delta_{\mathrm{KeyGen}}=0\)을 뜻하지 않는다. 실제 matrix helper의
\identifier{output_row}를 full-invNTT/pointwise-word 식으로 펼치는 마지막
rewrite와 actual harness의 두 active row consequence가 77번째 대상으로
컴파일됐다. 그러나 odd-root 직교성/full-NTT convolution과
\identifier{Rq.poly}-to-security-list adapter가 없다. 또한 sampler, retry,
packer, public API와 분포/종료를 포함하지 않는다. 따라서 KeyGen은
\identifier{STOP-KG-NTT}로 동결되고 활성 전선은 MINCORE Sign으로 이동했다.

\subsection{현재 주장 행렬}

\small
\begin{longtable}{L{0.24\textwidth}L{0.17\textwidth}L{0.51\textwidth}}
\caption{주요 주장과 현재의 정확한 강도}\label{tab:claim-matrix}\\
\toprule
주장 & 상태 & 현재 확보된 의미 / 남은 경계 \\
\midrule
\endfirsthead
\toprule
주장 & 상태 & 현재 확보된 의미 / 남은 경계 \\
\midrule
\endhead
KeyGen 패킹 접두부(packed prefix) & \statusproved & 종료한 mode-2 생성 KeyGen 반환값에서
\(sk[0..992)=vk[0..992)\); 종료와 분포는 별도 \\
KeyGen→메모리 도달 가능성(reachability) & \statusproved & 내부 반환값으로부터 구체적
로컬 Verify 메모리 witness 구성; Week~5 API 정리에는 사용하지 않는 역사적 seam \\
개별 public API key copy helper & \statusproved & 실제 생성 exporter/importer의
992/1408-byte 접두부 보존 부분정확성 \\
정준 주소 연결(canonical address bridge) & \statusproved & 정수 ABI 주소와 \identifier{W64}
exporter/hash 주소의 왕복 및 992/1408/1474-byte region 특수화 \\
실제 원시 KeyGen 내보내기(actual raw export) & \statusproved & 실제 caller가 종료하면 분리된 외부
VK/SK 영역의 첫 992바이트가 같음; retry 종료는 별도 \\
실제 원시 Sign 추적(actual raw trace) & \statusproved & 실제 caller의 반환값/memory를 보존하며
SK import와 hash pointer/length 관찰값을 노출 \\
실제 Verify 승인 경로(accepted path) & \statusproved & raw/cryptolab exact trace와
\(accept\Rightarrow tail\); 성공 시 실제 hash 입력 바인딩 \\
원시 보조절차 \(\mu\)(raw-helper \(\mu\)) & \statusproved & 실제 흡수(absorb),
Keccak, 마무리(finalize), 64/32바이트 짜내기(squeeze)의 관계 \\
생성 최상위 제어(generated top control) & \statusproved & 실제
\identifier{_sf_mu_rawpre}와 \identifier{__verify_hash_mu} 원시 분기(raw branch)를
직접 비교 \\
도전값 접미부(challenge suffix) & \statusproved & 위치 64 이후 실제 \(\mu_{32}\) absorb 결과
상태와 위치 일치 \\
승인 생성 어댑터(accepted generated adapters) & \statusproved & caller trace 사실이 실제 생성
hash-call \(\mu\) prefix와 위치-64 suffix helper 정리를 인스턴스화 \\
Sign 출력 프레임(Sign-output frame) & \statusproved & 실제 raw Sign의 1474-byte signature와
8-byte siglen 쓰기가 분리된 VK/SK/pre/message 구간을 보존하는 부분정확성 \\
구간 국소 생성 \(\mu\)(region-local generated \(\mu\)) & \statusproved & 전체 memory 등식 없이 실제
VK/pre/message read region만으로 생성 hash 반환값의 prefix 관계를 증명 \\
실제 Sign→Verify 순서(actual sequence) & \statusproved & 실제 두 caller와 trace 순차 실행의
결과/최종 memory 동치, accept-to-tail 및 hash descriptor 사후 바인딩 \\
승인 원시 API \(\mu\) 합성(accepted composition) & \statuspartial & trace의 직접
\identifier{observed_mu} 등식만 남음; frame/locality/control은 닫힘 \\
생성/재전달 운반(product/replay transport) & \statusdeferred & 생성 hash return의 pRHL 결과를
completed trace의 stored \identifier{observed_mu}로 옮길 정당한 규칙이 없음 \\
서명 접두부 코덱(signature prefix codec) & \statusproved & 실제 pack/unpack의 challenge 32-byte와
low 1024-byte round trip; canonical 입력 아래 부분정확성 \\
서명 접두부 프레임(signature prefix frame) & \statuspartial & unpack low-array tail은 보존되지만 pack
signature tail과 challenge-output tail은 열림 \\
HBZ 준비/적용 잎(prepare/apply leaf) & \statusproved & 실제 leaf가 canonical \(1024\)개 계수를
13-symbol alphabet으로 옮겼다가 복원하고 bounded array tail을 보존 \\
HBZ 실제 절차 경계(actual procedure boundary) & \statusproved & focused/full extraction과 실제
signature pack/unpack extraction의 HBZ full procedure가 exact equivalent; inverse는
별도 \\
HBZ 표 인증서(table certificate) & \statusproved & 실제 encoder/decoder/symbol 배열이 양의
빈도, 1024-slot partition 및 reciprocal arithmetic certificate를 만족 \\
순수 rANS 단일 단계(pure single step) & \statusproved & 수학적 \(D_s(E_s(x))=x\)와 허용 범위의
실제 reciprocal fast-step 일치; normalization byte와 loop는 제외 \\
순수 rANS 바이트 스택(byte stack) & \statusproved & 0/1/2-byte normalization readback,
little-endian state parse, \(xs/cuts/bytes\) construction과
\([2^{23},2^{31})\) state-bound 보존; generated loop는 제외 \\
실제 인코딩 접미부 복사(actual encoded-suffix copy) & \statusproved & 실제
\identifier{__copy_encoded_suffix}의 점별 slice equality와 unused-tail frame;
Week~13은 actual encoder success에서 입력 정수 범위와 W64 무랩을 유도해 핵심
합성에 사용 \\
실제 rANS 결합 제어(actual harness control) & \statusproved & 실제 encoder--copy--decoder 호출과
encoder 반환값에 따른 실패/성공 분기 및 decoder field 초기화; 이 기존 제어
정리 자체의 inverse 결론은 없음 \\
배열/목록 접미부 연결(array/list suffix bridge) & \statusproved & 정확한
1024원소 목록/접미부 점화식, 구간/프레임 연산 및 한 기호 추적 확장 \\
실제 인코더 단어 단계(actual encoder word step) & \statusproved & 실제 mode-2
표의 W32/W64 역수 계산이 순수 빠른 단계와 같음; 단어 수준 등식 \\
실제 내부 정규화(actual inner normalization) & \statusproved & 실제 생성 반복문이
순수 꼬리, 정확한 0/1/2바이트, 커서와 초기 접두부 프레임을 보존하고 실제
guard에서 정확히 종료하는 부분정확성 \\
최종 직렬화(final serialization) & \statusproved & 네 실제 패턴 저장이 누적
정규화 추적 앞에 리틀엔디언 상태를 붙이고 초기 접두부 프레임을 보존하도록
실제 생성 절차 출구에 합성됨 \\
실제 인코더 성공 크기(actual success size) & \statusproved & 실제
\identifier{_rans_encode}의 실패 분기를 보존하며 성공 시
\(0\le off\le1020\), \(4\le1024-off\le1024\); 종료/성공 도달은 제외 \\
실제 rANS 인코더 정제(actual encoder refinement) & \statusproved & 실제 반환
\(bad\)의 실패/성공 분기; 성공 시 전체 \identifier{trace_bytes}, 정확한
오프셋--길이 등식과 초기 접두부 프레임을 주는 성공 조건부 부분정확성 \\
실제 rANS 디코더 정제(actual decoder refinement) & \statusproved & 정준
1024기호의 정확 \identifier{trace_bytes}, 실제 표와 상태 입력 아래 종료한
\identifier{_rans_decode}가 \(bad=0\), 정확한 소비 크기, 전체 기호 복원과
출력 꼬리 프레임을 만족; 종료/손실없음은 별도 \\
실제 rANS 핵심 역함수(actual core inverse) & \statusproved & 실제 인코더 실패 시
디코더 미실행/초기 상태 보존, 성공 시 실제 복사와 디코더의 \(bad=0\), 정확한
소비 크기, 1024기호 복원 및 꼬리 프레임을 한 harness에서 증명; 성공 도달과
종료는 별도 \\
실제 전체 HBZ 부모(actual full-HBZ parent) & \statusproved & production
SignaturePack/Unpack 경계에서 zero-size 실패 보존 또는 nonzero-size 성공 시
1024계수 복원, 크기/두 꼬리 프레임과 prepare-symbol 추적을 주는 성공 조건부
부분정확성; 고정 zero 입력의 실패 배제는 별도 Week~15 정리, 종료는 별도 \\
고정 all-6/zero-HBZ 성공 사후조건 & \statusproved & 종료한 실제 all-6
인코더와 production zero-HBZ pack/unpack이 \(bad=0\), nonzero size, 정확한
trace/복원/프레임을 만족; 종료, losslessness와 일반 입력 성공은 별도 \\
KeyGen matrix/finalizer 스냅샷 & \statuspartial & 실제 두 helper의
matrix/NTT 표현, finalizer 의미, KG-2형 low/high 분해, 조정된 \(s_2\)와
스냅샷 mod-\(2q\) 영 등식은 증명됨; faithful KG-1/KG-3/전체 KG-4와 논문 식은
odd-root convolution 및 security-list adapter 부재로
\identifier{STOP-KG-NTT} \\
MINCORE Sign (활성) & \statusspecified & 실제
\identifier{_sf_round_challenge_mode2}, \identifier{_sf_z_check},
\identifier{_sf_hint_mode2}의 accepted-core 합성이 다음 단일 목표; 아직 headline
정리는 미컴파일 \\
MINCORE Verify/제한 합성 & \statusspecified & 실제 focused core 정리와
decoded-object 경계 합성의 목표만 고정됨; 아직 개별 정리가 컴파일되지 않음 \\
전체 서명 코덱(full signature codec) & \statusblocked & (h) rANS inverse,
metadata, padding, residual frame과 canonical parse가 필요하며 (h)는 Week~16
범위에서 \statusdeferred \\
정당한 Sign 출력 꼬리 도달(legitimate tail reach) & \statuspartial & prefix inverse만 닫혔고 full
codec의 HBZ 층까지 진전; (h), metadata/padding, norm, hint,
response/arithmetic가 필요한 별도
기능정확성 의무 \\
공용 문맥 \(\mu\)(public-context \(\mu\)) & \statusspecified & 문맥 길이 바이트,
문맥, Verify 상위 비트 분기(high-bit branch)의 일치 정리 없음 \\
전체 도전값 호출(full challenge call) & \statuspartial & \(\mu\) 접미부와 목록
합동성(list congruence)은 있으나 highbits/LSB, 패딩/도메인(padding/domain),
짜내기/표본기(squeeze/sampler)가 열림 \\
전체 KeyGen/Sign/Verify FC & \statusspecified & KeyGen 스냅샷 leaf는
부분적으로 진전했으나 세 API의 paper refinement 정리는 없음 \\
구현 EUF-CMA & \statusspecified & 모든 API loss, query accounting, 정확한
paper/ROM instance가 필요 \\
논문 보안 reduction & \statuspartial & 기존 조건부 chain은 있으나 정확한
transcript/distribution/forking 인스턴스가 열림 \\
\bottomrule
\end{longtable}
\normalsize

\subsection{가정, 계약, 증명 의무의 구분}

최종 보안 정리에 남길 수 있는 암호학적 가정은 구체적인 HAETAE-2 MLWE,
BST-MSIS, ROM 게임뿐이다. 반면 다음 항목은 암호학적 가정으로 숨기면 안 된다.

\begin{itemize}
  \item public buffer의 유효성, 크기, 분리, alias 조건;
  \item canonical 길이와 mode-2 상수;
  \item retry의 종료/losslessness와 RNG/random-tape 해석;
  \item FIPS202 바이트 transcript, padding, domain separation, endianness;
  \item hyperball sampler와 accepted-signature 분포;
  \item FFT 고정소수점 오차와 tied-minimum 정책.
\end{itemize}

이들은 API refinement, 확률 프로그램, 수치해석에서 각각 증명해야 할 계약이다.
현재 작성 이론에는 새 \identifier{axiom}이나 proof hole을 두어 이 간극을
덮지 않는다.
